\section{Fragility Score Evaluation}
\label{sec:fragility_statistical_evaluation}

The selected XGBoost predictions are converted into within-quarter percentile ranks, with higher Fragility Scores indicating greater predicted downside risk relative to other stocks in the same quarter. Table~\ref{tab:fragility-score-performance} summarizes the score's held-out ranking performance.

\begin{table}[!htbp]
\centering
\footnotesize
\caption{Held-out Fragility Score ranking performance.}
\label{tab:fragility-score-performance}
\begin{tabular}{lr}
\toprule
Statistic & Value \\
\midrule
Mean stock-level $\rho$: score vs. realized drawdown & 0.738 \\
Mean decile-level $\rho$: decile vs. realized drawdown & 1.000 \\
Mean realized-drawdown spread: Decile 10 minus Decile 1 & 41.6\% \\
\bottomrule
\end{tabular}
\end{table}


\noindent
At the stock level, the mean quarterly Spearman correlation between the Fragility Score and realized 63-day maximum drawdown is strong. At the decile level, mean realized drawdown increases monotonically with the Fragility Score in every test quarter, while the average realized-drawdown difference between Fragility Score Decile~10 and Decile~1 exceeds 40 percentage points. Figure~\ref{fig:fragility-score-deciles} and the supporting outcomes in Appendix~\ref{app:fragility_score_supporting_outcomes} show that higher score deciles also exhibit greater downside volatility, worse five-day losses, and weaker cumulative returns.
